\documentclass{article}
\usepackage{amsmath,amssymb}
\usepackage[margin=1in]{geometry}
\usepackage[UTF8]{ctex}
\begin{document}

\begin{equation}
\begin{aligned}
\max\quad & \sum_{t=1}^{T} f(\boldsymbol{x}_t;\boldsymbol{\theta}_t) \\
\text{s.t.}\quad & \sum_{t=1}^{T} g_i(\boldsymbol{x}_t;\boldsymbol{\theta}_t)\le c_i,
&& i=1,\ldots,m, \\
& \boldsymbol{x}_t\in\mathcal{X},
&& t=1,\ldots,T.
\end{aligned}
\tag{PCP}
\end{equation}

\paragraph{决策变量和固定数据（见~第7-8页）。}
决策变量为$(\boldsymbol{x}_1,\ldots,\boldsymbol{x}_T)$。时间范围由$T$个离散的时间段组成，而$i=1,\ldots,m$则用于索引预算限制。数量$c_i$表示预算$i$的初始容量。行动集满足……
\[
\mathcal{X}\subset\mathbb{R}^{k},\qquad
\mathcal{X}\text{ is compact and convex}.
\]
对于全时域计划而言，实际得到的参数属于数据范畴，而非决策变量。
\[
\boldsymbol{\theta}_t\in\Theta\subset\mathbb{R}^{l},
\qquad t=1,\ldots,T.
\]
函数$f$和$\boldsymbol{g}$是已知的。在这个选定的版本中，由于动作和参数的存在，这些函数具有确定性。而在第~6.4节中引入的奖励和消费方面的随机性在这里并未被使用。同时，也没有假设过时间范围$T$与初始预算之间的关系（p.~18）。

\paragraph{继承的函数类与参数化方式（pp.~7--8）。}
在底层定义中使用的空格保持不变。
\[
\begin{aligned}
\mathcal{F}=\mathcal{F}(\mathcal{X})
&=\{u:\mathcal{X}\to\mathbb{R}:u\text{ is concave and continuous on }\mathcal{X}\},\\
\mathcal{G}=\mathcal{G}(\mathcal{X})
&=\{u:\mathcal{X}\to\mathbb{R}:u\text{ is convex and continuous on }\mathcal{X}\},\\
f_t&\in\mathcal{F}, &&t=1,\ldots,T,\\
g_{it}&\in\mathcal{G}, &&i=1,\ldots,m,\quad t=1,\ldots,T.
\end{aligned}
\]
原始 p.~8 上显示的参数化信息如下：
\[
f_t(\boldsymbol{x}_t):=f(\boldsymbol{x}_t;\boldsymbol{\theta}_t),
\qquad
g_{it}(\boldsymbol{x}_t;\boldsymbol{\theta}_t):=g_i(\boldsymbol{x}_t;\boldsymbol{\theta}_t),
\qquad i=1,\ldots,m,\quad t=1,\ldots,T.
\]
在 (PCP) 中使用的向量值约束函数如下：
\[
\boldsymbol{g}(\boldsymbol{x};\boldsymbol{\theta})
=\bigl(g_1(\boldsymbol{x};\boldsymbol{\theta}),\ldots,
 g_m(\boldsymbol{x};\boldsymbol{\theta})\bigr)^{\top}
:\mathcal{X}\to\mathbb{R}^{m}.
\]
在 p.~7 上，继承下来的时间索引向量函数被以原始顺序打印出来。所有数字都保持不变。
\[
\boldsymbol{g}_t(\boldsymbol{x})
=\bigl(g_{1t}(\boldsymbol{x}),\ldots,g_{mt}(\boldsymbol{x})\bigr)^{\top}
:\mathbb{R}^{k}\to\mathbb{R}^{m}.
\]
\emph{符号表示上的差异被保留了下来：在}中，p.~8表示法在$g_{it}$的定义左侧包含了额外的参数参数；而在p.~7中，则直接引用了$g_{it}(\boldsymbol{x})$。此外，p.~7将$\boldsymbol{g}_t$的定义域表示为$\mathbb{R}^{k}$，尽管其组成部分的函数类是在$\mathcal{X}$上支持的。所选的p.~8表示法中的$\boldsymbol{g}$的定义域为$\mathcal{X}$；这里并没有提供该定义域之外的扩展情况。

\paragraph{参数空间与概率度量（pp.~8--9）。}
对于$\boldsymbol{\theta},\boldsymbol{\theta}'\in\Theta$，该论文中定义了：
\begin{equation}
\rho(\boldsymbol{\theta},\boldsymbol{\theta}')
:=\sup_{\boldsymbol{x}\in\mathcal{X}}
\left\|
\bigl(f(\boldsymbol{x};\boldsymbol{\theta}),\boldsymbol{g}(\boldsymbol{x};\boldsymbol{\theta})\bigr)
-
\bigl(f(\boldsymbol{x};\boldsymbol{\theta}'),\boldsymbol{g}(\boldsymbol{x};\boldsymbol{\theta}')\bigr)
\right\|_{\infty}.
\tag{2}
\end{equation}
这里的$\|\cdot\|_{\infty}$是$\mathbb{R}^{m+1}$中的 supremum norm。作者认为$\Theta$是一组类代表元素。
\[
\rho(\boldsymbol{\theta},\boldsymbol{\theta}')>0
\qquad\text{for every }\boldsymbol{\theta}\ne\boldsymbol{\theta}'\text{ in }\Theta.
\]
sigma-代数$\mathcal{B}_{\Theta}$是包含$\Theta$中所有在度量$\rho$下的开放子集的最小sigma-代数。分布$\mathcal{P}_t$是$\boldsymbol{\theta}_t$在$(\Theta,\mathcal{B}_{\Theta})$上的概率测度。

\paragraph{固定的输入假设（假设~1，页码~9）。}
这些是对于函数、动作集以及随机输入的限制，与(PCP)中的决策约束分开考虑。
\begin{align*}
\text{(a)}\quad
&|f(\boldsymbol{x};\boldsymbol{\theta})|\le 1
&&\forall\boldsymbol{x}\in\mathcal{X},\ \forall\boldsymbol{\theta}\in\Theta;\\
\text{(b)}\quad
&g_i(\boldsymbol{x};\boldsymbol{\theta})\in[0,1]
&&\forall\boldsymbol{x}\in\mathcal{X},\ \forall\boldsymbol{\theta}\in\Theta,
\quad i=1,\ldots,m;\\
&\boldsymbol{0}\in\mathcal{X},\qquad g_i(\boldsymbol{0};\boldsymbol{\theta})=0
&&\forall\boldsymbol{\theta}\in\Theta,\quad i=1,\ldots,m.
\end{align*}
\[
\text{(c)}\quad
\exists q>0\ \text{such that}\ 
\forall\boldsymbol{\theta}\in\Theta,\ \forall i=1,\ldots,m,
\ \forall\boldsymbol{x}\in\mathcal{X},\qquad
 g_i(\boldsymbol{x};\boldsymbol{\theta})>0
\ \Longrightarrow\ 
 f(\boldsymbol{x};\boldsymbol{\theta})\le q\cdot g_i(\boldsymbol{x};\boldsymbol{\theta}).
\]
\[
\text{(d)}\quad
\boldsymbol{\theta}_t\sim\mathcal{P}_t,\qquad t=1,\ldots,T.
\]
附带的印刷独立条件为：“$\mathcal{P}_t$的分布是彼此独立的。”这一表述将独立性直接赋予了分布本身；它被保留下来，而不是被另一个概率方程所替代。周围的随机输入描述使用跨周期的独立抽样，不要求分布完全相同，并且允许$(f,\boldsymbol{g})$的各个组成部分之间存在依赖关系。在一般模型中，并不假设对分布有精确的了解。
\[
\begin{aligned}
\text{(e)}\quad
&f(\boldsymbol{x};\boldsymbol{\theta})\text{ is concave over }\boldsymbol{x},\\
&g_i(\boldsymbol{x};\boldsymbol{\theta})\text{ is convex over }\boldsymbol{x},
\qquad\forall\boldsymbol{\theta}\in\Theta,\quad i=1,\ldots,m.
\end{aligned}
\]
条件(c)仅适用于所指示的消耗量严格为正的情况。特别地，所显示的假设并不要求在零动作时施加奖励归一化。不过，附录~A4中的内容（p.~47）仍然提到了在该动作下零目标值的情况；这一内容并未被添加为(PCP)的条件之一。

\paragraph{所选表格的范围为：pp.~7–8 和 11–12。}
该论文采用 (PCP) 作为其基本问题模型，即先观察当前状态，然后再做出决策。具体来说，当前参数、当前的奖励函数以及消费函数都是在选择当前行动之前被观察到的。所展示的完整视界程序是离线基准测试，而不是对在线策略的优化过程。因此，在线历史限制、后悔度量、期望价值松弛、双重问题以及算法子问题都不会为 (PCP) 增加新的行。

\end{document}
